\documentclass{article} % This command is used to set the type of
% document you are working on such as an article, book, or presenation
\usepackage{mathtools}
\usepackage{geometry} % This package allows the editing of the page layout
\usepackage{amsmath}  % This package allows the use of a large range
\usepackage{amsfonts}
% of mathematical formula, commands, and symbols
\usepackage{graphicx}  % This package allows the importing of images

\newcommand{\question}[2][]{
  \begin{flushleft}
    \textbf{Question #1}: \textit{#2}

\end{flushleft}}
\newcommand{\sol}{\textbf{Solution}:} %Use if you want a boldface solution line
\newcommand{\maketitletwo}[2][]{
  \begin{center}
    \Large{\textbf{Assignment #1}

    Topology} % Name of course here
    \vspace{5pt}

    \normalsize{N. Ohayon Rozanes  % Your name here

    \today}        % Change to due date if preferred
    \vspace{15pt}

\end{center}}
\begin{document}
\maketitletwo[11]  % Optional argument is assignment number
%Keep a blank space between maketitletwo and \question[1]

\question[1]{}
\begin{enumerate}
  \item
    \[
      (f \star e_q)(t)  =
      \begin{cases}
        f(2t) & 0 \le t \le \frac{1}{2}\\
        e_q(2t-1) & \frac{1}{2} \le t \le 1
      \end{cases} =
      \begin{cases}
        f(2t) & 0 \le t \le \frac{1}{2}\\
        q & \frac{1}{2} \le t \le 1
      \end{cases}
    \]
    \[
      (e_p \star f)(t)  =
      \begin{cases}
        e_p(2t) & 0 \le t \le \frac{1}{2}\\
        f(2t-1) & \frac{1}{2} \le t \le 1
      \end{cases} =
      \begin{cases}
        p & 0 \le t \le \frac{1}{2}\\
        f(2t-1) & \frac{1}{2} \le t \le 1
      \end{cases}
    \]
  \item One such homotopy from $f$ to $f \star e_q$.

    \[
      F(x,t) =
      \begin{cases}
        f(x(1-t) + t(2x)) & x \le \frac{1}{2} \\
        f(x(1-t) + t) & x \ge \frac{1}{2}
      \end{cases}
    \]
    To see that this is a homotopy, consider that:
    \[
      F(x, 0) = f(x) \quad F(x, 1) =
      \begin{cases}
        f(2x) & x\le \frac{1}{2} \\
        q & x\ge \frac{1}{2}
      \end{cases}
      = (f\star e_q)(x)
    \]
    Further, that $F$ is continuous is given by the fact that in $t$,
    the functions $t \to x(1-t)+ t(2x)$ and $t \to x(1-t) + t$ are
    straight lines in $\mathbb{R}^2$, and are therefore both
    continuous. Then since $f$ is also continuous, both $f \circ (t \to x(1-t)
    + t(2x))$ and $f \circ (t \to x(1-t) + t)$ are continuous. Then
    by the pasting lemma, $F(x,t)$ is a continuous function, and
    therefore is a homotopy from $f \to f\star e_q$

\end{enumerate}

\question[2]{}

\begin{enumerate}
  \item
    \[
      f \star \bar f =
      \begin{cases}
        f(2t)&  t \le \frac{1}{2} \\
        \bar f(2t - 1) & t \ge \frac{1}{2}
      \end{cases}
      =
      \begin{cases}
        f(2t)&  t \le \frac{1}{2} \\
        f(2(1 - t)) & t \ge \frac{1}{2}
      \end{cases}
    \]
    \[
      \bar f \star  f =
      \begin{cases}
        \bar f(2t)&  t \le \frac{1}{2} \\
        f(2t - 1) & t \ge \frac{1}{2}
      \end{cases}
      =
      \begin{cases}
        f(1-2t)&  t \le \frac{1}{2} \\
        f(2t-1) & t \ge \frac{1}{2}
      \end{cases}
    \]
  \item A homotopy from $e_p \to f \star \bar f$ is
    \[
      F(x,t) = f \star \bar f (tx)
    \]
    We have that
    $$F(x, 0) = f\star\bar f(0) = f(0) = p = e_p \quad F(x,1) = f(x)$$
    Since $f$ and therefore $\bar f$ are continuos, and therefore so
    is $f\star \bar f$ by the pasting lemma. Further, $t \to tx$ is
    clearly continuous as well, so $f \star \bar f \circ (t \to tx)$
    is continuous. As such $F(x,t)$ is continuous and therefore a homotopy.
\end{enumerate}

\question[3]{}
\begin{enumerate}
  \item \[
      f \star g =
      \begin{cases}
        f(2t) & t \le \frac{1}{2} \\
        g(2t - 1) & \frac{1}{2} \le t
      \end{cases}
    \]
    \[
      (f \star g) \star h =
      \begin{cases}
        f\star g (2k) & k \le \frac{1}{2} \\
        h(2k-1) & \frac{1}{2} \le k
      \end{cases} =
      \begin{cases}
        f(4t)  & t \le \frac{1}{4} \\
        g(4(t-\frac{1}{4}))  &  \frac{1}{4} \le t \le \frac{1}{2} \\
        h(2t-1) & \frac{1}{2} \le t
      \end{cases}
    \]
    similarily:

    \[
      f\star (g \star h) =
      \begin{cases}
        f(2t) & 0 \le t \le \frac{1}{2} \\
        g(4(t - \frac{1}{2})) & \frac{1}{2} \le t \le \frac{3}{4} \\
        h(4(t- \frac{3}{4})) & \frac{3}{4} \le t \le 1
      \end{cases}
    \]
  \item A homotopy from $(f \star g) \star h \to f \star (g \star h)$
    is a series of straight line homotopies.
\end{enumerate}
\question[4]{}
\[
  f \star g =
  \begin{cases}
    f(2t) & t \le \frac{1}{2} \\
    g(2t - 1) & \frac{1}{2} \le t
  \end{cases} \quad
  h \circ (f \star g) =
  \begin{cases}
    (h \circ f)(2t) & t \le \frac{1}{2} \\
    (h \circ g)(2t - 1) & \frac{1}{2} \le t
  \end{cases}
\]
And also

\[
  (h \circ f) \star (h \circ g) =
  \begin{cases}
    h\circ f (2t) & t \le \frac{1}{2} \\
    h \circ g (2t - 1) & \frac{1}{2} \le t
  \end{cases}
\]

Then clearly $(h \circ f) \star (h \circ g) =  h \circ (f \star g)$
\question[5]{}
First consider the line between $h:= h(x)$ and $x$, given by
\[
  h + v = x \to v = x - h
\]

Then the vector $v$ scaled by some constant $k \in \mathbb{R}$ from
$h$ will eventually intersect the unit circle, that point will be $r(x)$

\[
  kv + h = r(x) \quad ||r(x)||^2 = 1
\]
Then
\begin{align*}
  ||kv + h||^2 &= 1 \\
  ||k(x-h) + h||^2 &= 1 \\
\end{align*}
Expanding the above to vector components
\[
  k^2((x_1 - h_1)^2 + (x_2 - h_2)^2)  -  2k(h_1(x_1-h_1) + h_2(x_2 -
  h_2)) + (h_1^2 + h_2^2)-1= 0
\]
Then collecting this back into vector form, we get the follwing quadratic in $k$
\[
  k^2||v||^2 - 2k(h\cdot v) + (||h||^2 - 1) = 0
\]
Then through the quadratic equation we get that

\[
  k = \frac{h\cdot v + \sqrt{||v||^2 - ||h \times v||^2}}{||v||^2}
\]

With that, the closed form solution for $r$ is

$$
r(x) =  \frac{h\cdot v + \sqrt{||v||^2 - ||h \times v||^2}}{||v||^2} (x - h) + h
$$

\end{document}
